本文将避免医疗需求过载作为现实动机，但直接控制的是社区非隔离感染人数。第\ref{sec:model:strategy}节给出的联系经过全部新增感染流量，因此覆盖后来进入 $I_q$ 的病例，而不是把隔离病例排除在医疗需求之外。这与 Zhou 等\cite{Zhou2024SPCI}在完整模型中保留两类病例住院流入的处理相容；该文支持控制非隔离感染峰值的传播学理由，并未替本文建立容量充分条件。Zhang 等\cite[注记1]{Zhang2026Behavior}的感染--住院比例是其模型的平衡态关系，也不作为本文时变轨迹的固定换算。

上述平均需求界依赖共同的需求函数以及初始病例未来需求的上界，采用 $c_0S_0$ 得到对全部允许分配统一的充分条件，可能偏保守。它不要求隔离与非隔离感染人数成固定比例，但假定两类病例及不同分配采用相同的感染后需求关系；若医疗需求随病例来源、感染时点或干预措施改变，则需要重新建立共同上界。满足条件表示所设平均需求不超过容量，不表示真实随机占用的超载概率为零，也不能据此量化死亡率变化。第\ref{sec:xian:threshold-scale}节及联合控制比较中的 $\theta=0.002$ 仍是设计情景，未完成这一容量校准。$I+I_q$ 的峰值只反映总现存感染人数，不直接等同于住院或 ICU 占用。
